\section{Metodología y baseline}
\subsection{Preguntas de investigación M1/M2/M3}
Las preguntas de lectura se refieren a rendimiento, recursos y
comportamiento bajo concurrencia. La identificación normativa M1/M2/M3
debe conservar la formulación de los playbooks. Los extractos reconocidos
se presentan sin inventar una correspondencia.
\input{\artifactdir/preguntas.tex}

\subsection{Generación}
El responsable confirma 500\,000 documentos, semilla 42 y
1\,499\,443 ítems, aproximadamente tres por orden. La estructura incluye
órdenes, clientes, ubicaciones, productos, pagos y metadatos.
El desanidado produce una cardinalidad intermedia distinta del número
de documentos.

La semilla favorece repetibilidad, pero también importan generador y
dependencias. Por ello se conservan manifiesto y hashes. Las distribuciones
no se deducen de una descripción general.

SHA-256 comunicado del JSONL:
\begin{quote}\small
\nolinkurl{9fecf656a02cf6d369517ae4d59f47072e4de4025ae9db6280611347fb557afb}
\end{quote}
Hash lógico comunicado:
\begin{quote}\small
\nolinkurl{3490c8c2b9a58f58ee69849dd55f00d556e64456168964d0cf74262e95cc7715}
\end{quote}
\input{\artifactdir/dataset.tex}

\subsection{Carga y equivalencia}
Se comunican 500\,000 documentos importados en cada motor y equivalencia
lógica completa PASS. No implica igualdad binaria entre JSONL, BSON y
JSONB ni igualdad de espacio físico.

La igualdad de documentos no sustituye validar resultados de Q1/Q2.
Tipos, fechas, orden y empates requieren los controles del proyecto.
Importación y equivalencia son precondiciones, no resultados de rendimiento.

\subsection{Índices y aceptación}
MongoDB dispone, según el responsable, de \texttt{\_id} y de índices
simples en \texttt{timestamp} y \texttt{order\_id}.
PostgreSQL tiene PRIMARY KEY en \texttt{id} y B-Tree textual en
\texttt{data\textgreater\textgreater{}'timestamp'}.

Q2 convierte ese texto a \texttt{timestamptz}. Se distinguen:
\begin{enumerate}
\item existencia, definición y validez física;
\item compatibilidad con el predicado;
\item utilización efectiva en un plan;
\item aprobación académica.
\end{enumerate}
La expresión consultada importa para un índice sobre expresiones
\cite{pgindex}. La aceptación del responsable no demuestra uso,
compatibilidad con el cast ni aprobación del profesor.
La limitación permanece aunque las mediciones actuales sean Q1.

\subsection{Consultas}
Q1 desanida ítems, filtra \texttt{price > 50}, calcula subtotales y
agrupa por país y categoría, produciendo sumatoria, promedio y conteo.
Deben mantenerse expresiones y tipos originales.

Q2 selecciona órdenes completadas en un rango temporal, ordena por
importe final descendente y recupera Top-1000. No hay resultados
definitivos Q2 incorporados.

Los extractos proceden de archivos reales y se conservan fuentes completas.
La selección automática no acredita por sí sola su invocación por el runner.
\input{\artifactdir/consultas.tex}

\subsection{Variantes, matriz y protocolo}
La matriz base selecciona \texttt{allow\_disk\_false} y
\texttt{work\_mem\_4mb}. Las variantes \texttt{allow\_disk\_true}
y \texttt{work\_mem\_64mb} se mantienen separadas. No son presupuestos
internos equivalentes de memoria.

El diseño comprende:
\[
3\text{ perfiles}\times2\text{ motores}\times2\text{ consultas}
\times3\text{ concurrencias}\times3\text{ repeticiones}=108.
\]
Las concurrencias son 1, 10 y 50. Una repetición por escenario comprende
36 ejecuciones y no completa el protocolo.

Se establecen 60 s de warmup, 300 s de medición y 60 s de cooldown.
La duración total incluye preparación y controles: los 53 minutos y
38 segundos comunicados para PostgreSQL no sustituyen la ventana ni
se usan como denominador de QPS.

Debe conservarse el criterio del runner sobre operaciones en los límites
de ventana. No se reconstruye la latencia media desde la inversa de QPS
ni se inventa conteo de operaciones.

\subsection{Cachés y secuencia}
El warmup no acredita estabilización. Caché del motor, caché del sistema
operativo y conexiones son controles distintos. Un comando en un script
no demuestra que se ejecutara en cada repetición.

Este generador no limpia cachés, reinicia servicios ni ejecuta operaciones
privilegiadas. Las condiciones deben cotejarse con evidencia existente.

\subsection{Preflight e incidencia}
El responsable confirma 135 pruebas automatizadas aprobadas para la
optimización. La reutilización se bloqueó al cambiar de motor.
Los casos probados y el intento de reutilización son evidencias distintas.

No hay comparación temporal válida antes/después. No se afirma ahorro
experimental. Se registra la incidencia sin modificar resultados
comprometidos ni suprimir controles.

\subsection{Contrato}
La cabecera conserva 14 columnas:
\begin{lstlisting}
timestamp_epoch,engine,profile,query_id,concurrency,run_number,throughput_qps,latency_p50_ms,latency_p95_ms,latency_p99_ms,cpu_throttled_seconds,max_rss_mb,disk_spill,status
\end{lstlisting}
\texttt{spill} puede coexistir con elegibilidad. Las ausencias no se
completan con cero. La confirmación del responsable se distingue del
reconocimiento automático de campos en las evidencias.
